\section*{अध्याय \ref{ch:b3:exoplanets} --- बहिर्ग्रह}

\begin{solution}{exo:b3:exoplanets:1}
(क)~$K = (2\pi G/P)^{1/3}m_p/\bigl((M_\odot + m_p)^{2/3}\sqrt{1 - e^2}\bigr)$।
बृहस्पति के लिए: $P = \qty{3.742e8}{s}$, $(2\pi G/P)^{1/3} = \num{1.039e-6}$ SI
इकाइयों में, $K = \num{1.039e-6}\times\num{1.898e27}/(\num{1.990e30})^{2/3}/0.9988
= \qty{12.5}{m.s^{-1}}$। पृथ्वी के लिए: $P = \qty{3.155e7}{s}$, $K =
\qty{0.089}{m.s^{-1}}$, लगभग \qty{9}{cm.s^{-1}}। (ख)~\qty{3.5}{km/h} $=
\qty{0.97}{m.s^{-1}}$ और \qty{0.35}{km/h} $= \qty{0.097}{m.s^{-1}}$।
बृहस्पति का संकेत पहली परिशुद्धता का तेरह गुना है; पृथ्वी का
दूसरी परिशुद्धता से भी नीचे है, और तारे का अपना शोर उस पर और जुड़ जाता है।
\end{solution}

\begin{solution}{exo:b3:exoplanets:2}
(क)~$(71492/695700)^2 = \num{0.0106}$, \qty{1.06}{\percent} की गहराई;
$(6378/695700)^2 = \num{8.4e-5}$, 84 प्रति दस लाख भाग। (ख)~$R_\star/a$:
पृथ्वी के लिए $695700/\num{1.496e8} = \num{4.7e-3}$, लगभग 215 में से एक अवसर;
बृहस्पति के लिए $695700/\num{7.785e8} = \num{8.9e-4}$, 1100 में से एक अवसर।
\end{solution}

\begin{solution}{exo:b3:exoplanets:3}
$m_p\sin i \approx K M_\odot^{2/3}(P/2\pi G)^{1/3}$ जहाँ $P =
\qty{3.655e5}{s}$: $(P/2\pi G)^{1/3} = \num{9.55e4}$ SI इकाइयों में, अतः $m_p
\sin i = 59\times(\num{1.9884e30})^{2/3}\times\num{9.55e4} =
\qty{8.9e26}{kg}$, अर्थात 0.47 बृहस्पति-द्रव्यमान या 149 पृथ्वी-द्रव्यमान (
$m_p$ के लिए सुधार, $M_\star + m_p$ में, हज़ार में एक भाग से कम है)।
$a = \bigl(GM_\odot P^2/4\pi^2\bigr)^{1/3} = \qty{7.66e9}{m} =
\qty{0.051}{au}$।
\end{solution}

\begin{solution}{exo:b3:exoplanets:4}
(क)~$4GM/c^2 = 4\times\num{6.674e-11}\times\num{9.94e29}/(\num{2.998e8})^2 =
\qty{2953}{m}$; $(D_S - D_L)/(D_LD_S) = 1/D_S\times(D_S - D_L)/D_L =
1/(\qty{8}{kpc}) = 1/\qty{2.469e20}{m}$। $\theta_\mathrm E =
\sqrt{2953/\num{2.469e20}} = \qty{3.46e-9}{rad} = \qty{0.71}{mas}$।
लेंस पर, $R_\mathrm E = \theta_\mathrm E D_L = \num{3.46e-9}\times
\qty{1.234e20}{m} = \qty{4.27e11}{m} = \qty{2.85}{au}$। (ख)~$A(1) =
3/\sqrt5 = 1.34$; $A(0.5) = 2.25/(0.5\times2.062) = 2.18$; $A(0.1) =
2.01/(0.1\times2.002) = 10.0$। (ग)~$t_\mathrm E = \num{4.27e11}/\num{2e5}
= \qty{2.13e6}{s}$, लगभग 25 दिन: \omterm{def:b3:exoplanets:microlensing}{गुरुत्वीय सूक्ष्म-लेंसन} की घटनाएँ सप्ताहों तक रहती हैं।
\end{solution}

\begin{solution}{exo:b3:exoplanets:5}
(क)~$a = (GM_\odot P^2/4\pi^2)^{1/3} = \qty{6.09e9}{m}$ (\qty{0.0407}{au}),
अतः $a/R_\star = 8.75$। (ख)~$P/\pi = \qty{22.9}{h}$ लेते हुए, $b = 0$ के लिए:
\[
T_{14} = \frac P\pi\arcsin\frac{1.12}{8.75} = \qty{2.94}{h}, \qquad
T_{23} = \frac P\pi\arcsin\frac{0.88}{8.75} = \qty{2.31}{h}.
\]
$b = 0.7$ के लिए: $\cos i = 0.7/8.75$, $i = \ang{85.4}$, $(a/R_\star)\sin i =
8.726$, और
\[
T_{14} = \frac P\pi\arcsin\frac{\sqrt{1.2544 - 0.49}}{8.726} = \qty{2.30}{h},
\qquad
T_{23} = \frac P\pi\arcsin\frac{\sqrt{0.7744 - 0.49}}{8.726} = \qty{1.40}{h}.
\]
(ग)~प्रवेश और निर्गम पारगमन का $(T_{14} - T_{23})/T_{14} = \qty{21}{\percent}$
भाग लेते हैं जब यह केंद्रीय हो, और $b = 0.7$ पर \qty{39}{\percent}; एक
स्पर्शी पारगमन, $b \to 1$, में उसमें कोई समतल खंड नहीं होता और यह V जैसा दिखता है।
\end{solution}

\begin{solution}{exo:b3:exoplanets:6}
$\bar\rho = M/\tfrac43\pi R^3$, जहाँ $\tfrac43\pi R_\oplus^3 =
\qty{1.087e21}{m^3}$: ट्रैपिस्ट-1~b $1.374\times\num{5.97e24}/(1.116^3\times
\num{1.087e21}) = \qty{5.43}{g.cm^{-3}}$; केप्लर-10~b
\qty{5.60}{g.cm^{-3}}; WASP-39~b \qty{0.167}{g.cm^{-3}}; HD~209458~b
\qty{0.337}{g.cm^{-3}}। पहले दो पृथ्वी जितने ही घने हैं
(\qty{5.49}{g.cm^{-3}}, वही स्थिरांक लेकर) और चट्टानी; अंतिम दो
गैस दानव हैं, पानी से हल्के; WASP-39~b
बृहस्पति से आठ गुना हल्का है।
\end{solution}

\begin{solution}{exo:b3:exoplanets:7}
चरम मान $\gamma + K(1 + e\cos\omega)$ और $\gamma - K(1 -
e\cos\omega)$ हैं। इनका अंतर $2K = \qty{60.0}{m.s^{-1}}$ है, अतः $K =
\qty{30.0}{m.s^{-1}}$; इनका माध्य $\gamma + Ke\cos\omega =
\qty{20.4}{m.s^{-1}}$ है, अतः $Ke\cos\omega = \qty{8.0}{m.s^{-1}}$ और
$e\cos\omega = 0.27$। उत्केंद्र कक्षा में तारा अपसौर के पास ठहरता है
और उपसौर से तेज़ी से गुज़रता है, इसलिए वक्र असममित होता है और
इसके चरम मानों का मध्य-बिंदु उसका समय-औसत $\gamma$ नहीं है।
\end{solution}

\begin{solution}{exo:b3:exoplanets:8}
(क)~$g = GM/R^2 = \num{6.674e-11}\times\num{5.332e26}/(\num{9.144e7})^2 =
\qty{4.26}{m.s^{-2}}$; $H = kT/\bar m g = \num{1.381e-23}\times1100/(2.3
\times\num{1.661e-27}\times4.26) = \qty{934}{km}$। $R_\star =
\qty{6.533e8}{m}$ लेने पर, गहराई $(\num{9.144e7}/\num{6.533e8})^2 =
\num{0.0196}$ होती है। (ख)~$2HR_p/R_\star^2 = 2\times\num{9.34e5}\times
\num{9.144e7}/(\num{6.533e8})^2 = \num{4.0e-4}$, 400 प्रति दस लाख भाग;
दो पैमाना-ऊँचाइयाँ, यानी 800; कम गुरुत्व वाला गैस दानव
\omterm{def:b3:exoplanets:transmission}{पारगमन स्पेक्ट्रमिकी} के लिए सबसे आसान लक्ष्य होता है।
\end{solution}

\begin{solution}{exo:b3:exoplanets:9}
$p = 1.116\times6378/(0.1192\times695700) = 0.0858$, $p^2 = \num{7.37e-3}$।
$B_\lambda \propto 1/(e^{hc/\lambda kT} - 1)$ जहाँ $hc/\lambda k =
\qty{959}{K}$, \qty{15}{\micro m} पर: तारा $1/(e^{0.374} - 1) = 2.20$;
दिवस-पक्ष \qty{400}{K} पर, $1/(e^{2.40} - 1) = 0.0998$, अनुपात \num{0.0453},
गहराई \num{3.3e-4} (334 प्रति दस लाख भाग); \qty{500}{K} पर, $1/(e^{1.92} -
1) = 0.172$, अनुपात \num{0.078}, गहराई \num{5.7e-4}। सौ केल्विन
का अंतर गहराई को \qty{70}{\percent} बदल देता है: मापी गई गहराई दिवस-पक्ष
के तापमान को निकटता से निर्धारित करती है, और गर्म दिवस-पक्ष का अर्थ है कि
बहुत कम ऊष्मा रात्रि-पक्ष तक ले जाई जाती है।
\end{solution}

\begin{solution}{exo:b3:exoplanets:10}
(क)~$S \propto L/d^2$, अतः $d = \sqrt{0.005}\,\unit{au} = \qty{0.0707}{au}$;
$P = 2\pi\sqrt{a^3/GM} = \qty{1.33e6}{s} = \qty{15.4}{d}$।
(ख)~$K = (2\pi G/P)^{1/3}M_\oplus/(0.2M_\odot)^{2/3} = \qty{0.75}{m.s^{-1}}$,
सूर्य पर पृथ्वी के प्रभाव का 8.4 गुना: हल्का तारा अधिक गति करता है, और
छोटा आवर्तकाल भी सहायक है। (ग)~$(6378/139140)^2 = \num{2.1e-3}$, पृथ्वी के
सूर्य-पारगमन से 25 गुना। दोनों संकेत बड़े हैं, पारगमन
हर दो सप्ताह में आते हैं, और पारगमन प्रायिकता $R_\star/a =
\num{0.013}$ पृथ्वी की तुलना में तीन गुना है: छोटे तारे
उनके वासयोग्य क्षेत्र में चट्टानी ग्रहों को खोजे और अध्ययन किए जाने का अवसर देते हैं।
\end{solution}

\begin{solution}{exo:b3:exoplanets:11}
$r\,dr = -\mu\,d\mu$ के साथ, माध्य तीव्रता $\frac1\pi\int_0^1
I\,2\pi r\,dr = 2\int_0^1 I(\mu)\mu\,d\mu = 2\bigl[\tfrac12 - u_1(\tfrac12 -
\tfrac13) - u_2\int_0^1(1 - \mu)^2\mu\,d\mu\bigr]$ अंतिम समाकल
$\tfrac12 - \tfrac23 + \tfrac14 = \tfrac1{12}$ है, अतः माध्य $1 - u_1/3 -
u_2/6$, यहाँ $0.825$। एक छोटा ग्रह $\pi p^2 I(\mu)$ ढकता है, $\pi\times
0.825$ का भाग: गहराई $p^2 I(\mu)/0.825$। $r = 0$ पर, $I = 1$: $1.212\,p^2$;
$r = 0.5$ ($\mu = 0.866$) पर, $I = 0.942$: $1.142\,p^2$; $r = 0.9$ ($\mu =
0.436$) पर, $I = 0.695$: $0.842\,p^2$। एक समान-चकती मॉडल
$1.212\,p^2$ को $p^2$ पढ़ेगा और त्रिज्या को $\sqrt{1.212} = 1.10$ से
अधिक आँकेगा, \qty{10}{\percent}।
\end{solution}

\begin{solution}{exo:b3:exoplanets:12}
(क)~$R_\odot/a$: \num{9.3e-3} \qty{0.5}{au} पर, \num{6.6e-3}
\qty{0.7}{au} पर, \num{4.7e-3} \qty{1}{au} पर। (ख)~$42000\times0.11\times
\num{6.6e-3} = 31$ ग्रह पारगमन करते हैं। (ग)~इनमें से लगभग एक-तिहाई ही
मिले: बाकी के पारगमन इतने उथले थे (पृथ्वी के लिए 84 पार्ट्स पर मिलियन
) या इतने कम (प्रति वर्ष एक या दो) कि शोर से ऊपर न उठ सके। वह
अंश संसूचन दक्षता है; इसके लिए और पारगमन प्रायिकता के लिए सुधार किए बिना,
गणना ऐसे ग्रहों की सामान्यता
को लगभग 460 गुना कम आँकती।
\end{solution}

\begin{solution}{pb:b3:exoplanets:1}
\textbf{1.} $M_\star = 0.80\times\num{1.9884e30} = \qty{1.591e30}{kg}$, $P =
\qty{1.0368e6}{s}$: $a = (GM_\star P^2/4\pi^2)^{1/3} = \qty{1.4245e10}{m}
= \qty{0.0952}{au}$।

\textbf{2.} $v = 2\pi a/P = \qty{86.3}{km.s^{-1}}$।

\textbf{3.} $m_p\sin i = KM_\star^{2/3}(P/2\pi G)^{1/3} = 2.59\times\num{1.363e20}
\times\num{1.352e5} = \qty{4.77e25}{kg}$, \num{7.99} पृथ्वी द्रव्यमान।

\textbf{4.} $a_\star = a\,m_p/M_\star = \num{1.4245e10}\times\num{4.77e25}/
\num{1.591e30} = \qty{427}{km}$; $2\pi a_\star/P = \qty{2.59}{m.s^{-1}}$,
देखा गया $K$।

\textbf{5.} $2.59$ गुना \qty{1}{m.s^{-1}} परिशुद्धता का, 26 गुना एक
\qty{0.1}{m.s^{-1}} की: ग्रह सर्वोत्तम उपकरणों की पहुँच में है,
पिछली पीढ़ी की पहुँच में मुश्किल से।

\textbf{6.} $m_p\sin i \propto K\sqrt{1 - e^2}$: $\times0.980$, 7.83 पृथ्वी
द्रव्यमान, \qty{2}{\percent} कम।

\textbf{7.} $p = \sqrt{\num{8.38e-4}} = 0.02895$।

\textbf{8.} $R_\star = 0.76\times\qty{695700}{km} = \qty{528700}{km}$; $R_p =
\qty{15310}{km} = 2.40$ पृथ्वी त्रिज्याएँ।

\textbf{9.} $a/R_\star = \num{1.4245e10}/\num{5.287e8} = 26.94$।

\textbf{10.} $T_{14} = (\qty{288}{h}/\pi)\arcsin(1.02895/26.94) = \qty{3.50}{h}$।

\textbf{11.} $\sqrt{(1 + p)^2 - b^2} = (a/R_\star)\sin(\pi T_{14}/P) = 26.94\times
\sin(\pi\times3.23/288) = 0.949$ (जहाँ $\sin i \approx 1$); $b^2 =
1.0587 - 0.9008 = 0.158$, $b = 0.40$।

\textbf{12.} $\cos i = b\,R_\star/a = 0.40/26.94 = 0.0148$, $i = \ang{89.15}$।

\textbf{13.} $R_\star/a = 1/26.94 = 0.037$: 27 में से एक अवसर।

\textbf{14.} $\sin i = 0.99989$: वास्तविक द्रव्यमान 7.99 पृथ्वी द्रव्यमान है, \omterm{def:b3:exoplanets:rv}{न्यूनतम
द्रव्यमान} \num{1e-4} के भीतर है।

\textbf{15.} $\bar\rho = \num{4.77e25}/\bigl(\tfrac43\pi(\num{1.531e7})^3\bigr) =
\num{4.77e25}/\num{1.502e22} = \qty{3180}{kg.m^{-3}}$,
\qty{3.18}{g.cm^{-3}}।

\textbf{16.} ग्रह अपने द्रव्यमान की चट्टानी ग्रह से $2.40/1.74 = 1.38$ गुना बड़ा है,
आयतन में 2.6 गुना। केवल चट्टान से यह नहीं बन सकता: इसे
हल्की गैस का एक आवरण चाहिए, अपने द्रव्यमान का कुछ प्रतिशत हाइड्रोजन और हीलियम,
या जल का बड़ा अंश।

\textbf{17.} 2.40 पृथ्वी त्रिज्याओं पर यह \omterm{def:b3:exoplanets:radius-valley}{त्रिज्या घाटी} के ऊपर है, 2--3 पृथ्वी-त्रिज्या
वाले शिखर में: एक \omterm{def:b3:exoplanets:classes}{उप-नेपच्यून}।

\textbf{18.} $g = Gm_p/R_p^2 = \num{6.674e-11}\times\num{4.77e25}/(\num{1.531e7})^2 =
\qty{13.6}{m.s^{-2}}$।

\textbf{19.} $L = 4\pi R_\star^2\sigma T^4 = 4\pi(\num{5.287e8})^2\times\num{5.670e-8}
\times5000^4 = \qty{1.245e26}{W}$, 0.325 सौर दीप्तियाँ।

\textbf{20.} $S/S_\oplus = 0.325/0.0952^2 = 35.9$: वासयोग्य क्षेत्र की भीतरी
सीमा से बहुत अंदर।

\textbf{21.} $L/4\pi a^2 = \qty{4.88e4}{W.m^{-2}}$; $T_\mathrm{eq} =
\bigl(\num{4.88e4}\times0.7/(4\sigma)\bigr)^{1/4} = \qty{623}{K}$।

\textbf{22.} $H = kT/\bar m g = \num{1.381e-23}\times623/(2.3\times\num{1.661e-27}
\times13.6) = \qty{166}{km}$।

\textbf{23.} $2HR_p/R_\star^2 = 2\times\num{1.66e5}\times\num{1.531e7}/
(\num{5.287e8})^2 = \num{1.82e-5}$, 18 पार्ट्स पर मिलियन; दो
पैमाना ऊँचाई तक का पट्टा, 36 पार्ट्स पर मिलियन।

\textbf{24.} $H$ का मान $18/2.3 = 7.8$ गुना छोटा है, \qty{21}{km}: 2.3 पार्ट्स पर
मिलियन प्रति पैमाना ऊँचाई, पट्टे के लिए 4.6। दस पार्ट्स पर मिलियन पर,
हाइड्रोजन-समृद्ध वायुमंडल अपना जल-पट्टा शोर के तीन गुने से अधिक
परिमाण पर दिखाता है; भाप का वायुमंडल सपाट दिखता है, जैसे बादलों वाला
या बिना वायु का ग्रह।

\textbf{25.} \textbf{\qty[text-series-to-math=true]{3.2}{g/cm^3}}।
\end{solution}
